\chapter*{Colophon} 
\thispagestyle{empty}
\shapepar{\hexagonshape}{%(vedi computer science per layout colophon)
Ce livre a été écrit avec \LaTeX\ (La(mport)\TeX), un logiciel gratuit de composition typographique créé par \mbox{Leslie} Lamport, qui utilise \TeX, une invention de \mbox{Donald} Ervin Knuth datant de 1982, comme moteur de composition.
La \textbf{police} utilisée est \mbox{Charter}, une police de caractères à empattements (serif) conçue par le renommé typographe Matthew \mbox{Carter} en 1987. %Caractérisée par des formes claires et des contreformes ouvertes, elle est reconnue pour son excellente lisibilité, ce qui en fait un choix apprécié pour le corps de texte.
Les textes  en francoprovençal ont été révisés par le \mbox{\textbf{Guichet linguistique}} de l'Assessorat régional des activités et des biens culturels, du système educatif et des politiques des relations intergénérationnelles.
Les \textbf{émoticônes} \ok\ ont été aimablement fournies par Icons8 (\href{https://icons8.com/icons/}{https://icons8.com/icons/}).
La \mbox{\textbf{couverture}} du livre (\macchinaFoto\ Roger Berthod) reflète parfaitement l'esprit des Digourdì : ce désir de regarder au-delà du rideau, avec l'élan et la curiosité de la jeunesse. Mais aussi avec des coulisses dignes d’une véritable Comédie, qui nous apprend à vivre tels que nous sommes, tout simplement, avec nos imperfections.
}
% prendere spunto da colophon della dispensa Pantieri o computer science
\vfill 
\noindent\makebox[\textwidth][c]{%
  \busta\ \href{mailto:jbollon94@gmail.com}{jbollon94@gmail.com}
}
\noindent\makebox[\textwidth][c]{%
  \busta\ \href{mailto:digourdi.tsarvensou@gmail.com}{digourdi.tsarvensou@gmail.com}
}
\thispagestyle{empty}